\section{Grammar discovery and its finite limits}
\label{sec:grammar-search}

The semantic freedoms above were integrated into actual grammar search
programs.  None of the recorded campaigns produced a grammar with 127
roles, nor a 128-role grammar with maximum zero run five.  The positive
results of this part are a portable search representation, separately
checked finite candidate grammars, and an exact obstruction showing why
one apparently natural search family could never satisfy its objective.

\subsection{The representation and objective}

A grammar role has two outgoing decoder choices, represented by a rule
$w\to(u,v)$.  In the ordinary domain, $w=uv$ is literal concatenation of
inherited semantic symbols.  In the fine domain, every symbol is expanded
over the 30 closed atoms.  The 44 inherited implementations remain
available as zero-additional-cost terminals, including their 14 composite
expansions.  This exposes reassociations across the old symbolic
boundaries without charging again for already present implementation
roles.  Terminal implementations retain their actual allowed left/right
entry orientations.  The tokens in this representation are compiler
identifiers; the target tape alphabet remains binary.

In a trace domain, the concatenation condition becomes
$\operatorname{nf}(w)=\operatorname{nf}(uv)$, where $\operatorname{nf}$ is
the lexicographically least linearization of the dependence graph for
the specified independence relation.  The search can therefore share a
role between different certified orderings of the same closed
computation.  Its moves include binary reassociation, alternative cuts
using already available words, changes to folded-entry routes, and
construction of small router words.  The affine version also tries
commuting child swaps and adjacent commuting changes near a proposed cut.
No rule may refer cyclically to itself through another rule.

The objective counts all live grammar roles, including roles needed only
by the folded-entry routes.  Eleven routes satisfy 22 specified decoder
projection equations.  A candidate also retains required source entry
roots and, when requested, the 54 inherited bootstrap projection
equations.  These equations are sufficient inherited interfaces, rather
than a characterization of every possible valid initialization.

The scanner constraint is evaluated on the complete expanded source
tree, including fixed atom implementations and unchanged source
subtrees.  For a leaf put $(L,Z,D)=(0,0,0)$, where $L$ is the leading-zero
depth, $Z$ the maximum zero run, and $D$ the maximum leaf depth.  Binary
composition obeys
\[
 L(A,B)=1+L(A),\qquad
 Z(A,B)=\max\{1+L(A),Z(A),Z(B)\},\qquad
 D(A,B)=1+\max\{D(A),D(B)\}.
\]
The ordinary and trace searches retain $Z\leq6$ and $D\leq64$.  The
five-zero searches additionally minimize the number $B_6$ of tree-node
occurrences from which six consecutive left edges exist:
\[
 B_6(A,B)=B_6(A)+B_6(B)+\mathbf 1_{L(A)\geq5}.
\]
This counts occurrences in the expanded tree, including repetitions of a
shared grammar role.  It is not a count of reachable execution sites.

\subsection{Recorded annealing campaigns}

Table~\ref{tab:grammar-campaigns} reports the saved output of each
campaign.  Every listed output was checked again by the independent
Python grammar auditor.  It verifies concatenation in the stated domain,
the fixed terminal orientations, mandatory source roots, all 22 route
projections, the supplied bootstrap equations, and the expanded-tree
statistics.  The affine audit uses the separately justified 178-pair
relation.  It is a finite grammar audit, not a new blank-start or
whole-runtime proof.

\begin{table}[htbp]
\centering\small
\begin{tabularx}{\textwidth}{@{}Xrrrrr@{}}
\toprule
Campaign & Locks & Proposals & Roles & $Z$ & $B_6$ \\
\midrule
Ordered words & 54 & 7,099,557 & 128 & 6 & 49 \\
Ordered, unlocked & 0 & 3,938,413 & 128 & 6 & 49 \\
Fine words, warm & 54 & 3,937,175 & 128 & 6 & 48 \\
Fine words, cold & 54 & 2,497,065 & 128 & 6 & 49 \\
Disjoint-support traces & 54 & 863,648 & 128 & 6 & 50 \\
Faster trace search & 54 & 2,042,251 & 128 & 6 & 48 \\
Affine-extended traces & 54 & 5,073,185 & 128 & 6 & 50 \\
Five-zero, mild penalty & 50 & 6,670,998 & 128 & 6 & 44 \\
Five-zero, strong penalty & 50 & 4,070,474 & 134 & 6 & 6 \\
Five-zero, fully unlocked & 0 & 5,486,426 & 134 & 6 & 4 \\
\bottomrule
\end{tabularx}
\caption{Bounded discovery runs: 41,679,192 proposal iterations in total.
The last two outputs minimize a weighted score, so their reported role
counts can exceed the 128-role initial control.}
\label{tab:grammar-campaigns}
\end{table}

The mild five-zero score is $K+0.05B_6$, where $K$ counts live roles; the
strong score is $K+0.30B_6$.  Their uphill-proposal caps are respectively
133 and 137 roles.  The strong campaigns demonstrate that most six-zero
chains can be removed at a grammar cost, but both retain $Z=6$.  They do
not meet either numerical target.  The 50-lock runs use the association
rotation described below.  The fully unlocked run removes all bootstrap
equations and starts from the same rotated seed.  No modified blank
execution was certified for any released-lock output.

The proposals include rejected moves and repeated grammars.  The total is
not a count of distinct candidates and is not an exhaustive enumeration.
Each run is wall-time bounded, with seeds and settings in
\file{grammar/search_profiles.json}; runtime-dependent stopping and
implementation details can change a rerun's exact trajectory.  The saved
JSON and logs specify the actual investigated artifacts.  Their inherited
\texttt{framework\_states:118} metadata predates the existing scanner
contraction and must not be added uncritically to a grammar count to claim
a literal new machine.  These files describe abstract grammars only.

\subsection{Finite candidate optimization and reuse-biased parsing}

Three binary-production MILP models test increasingly large finite
dictionaries.  A word-activation variable chooses whether a word has a
role; a production variable chooses one legal split into available
children.  Every active word has exactly one active production, and all
nonterminal children of an active production must be active.  Required
roots are active.  Since children are shorter words, selected productions
are acyclic.  The objective is the number of active words, with the stated
upper bound.  Candidate pools contain the old words, inherited terminal
expansions, and substrings occurring at least twice among the selected
roots, through the stated length cap.  Unsupported words are pruned.

The fixed-route models retain 39 roots and 43 complete old productions,
the latter being enough to preserve the old bootstrap and route paths.
Thus they freeze more than the abstract route projection equations alone.
The dynamic-route model has 32 roots, 27 frozen productions, and 315
supported route options.  Its options are generated from the finite
production pool and explicitly enforce the required one- or two-bit
projections.  These models do not impose the scanner zero-run or depth
bounds.

\begin{table}[htbp]
\centering\small
\begin{tabularx}{\textwidth}{@{}lrrrX@{}}
\toprule
Model & Length cap & Words & Productions & Solver report \\
\midrule
Fixed routes & 16 & 1,004 & 5,059 & Infeasible at $K\leq127$ \\
Fixed control & 16 & 1,004 & 5,059 & Optimum $K=128$ at $K\leq128$ \\
Fixed routes & 64 & 2,720 & 31,572 & Infeasible at $K\leq127$ \\
Dynamic routes & 32 & 2,014 & 17,716 & Infeasible at $K\leq127$ \\
\bottomrule
\end{tabularx}
\caption{HiGHS reports about the specified finite MILP models.  No
independently checked infeasibility or optimality proof log is supplied.}
\label{tab:grammar-milp}
\end{table}

The 128-role control solution was independently checked as a valid finite
grammar.  This is a feasibility check of its explicit witness; it does not
independently establish the reported optimum.  The three infeasibility
reports concern only their generated finite pools and their stronger
frozen interfaces.  They are not lower bounds for the unrestricted
grammar problem or for Turing machines.

A further phrase-flow model lets an activated phrase be decomposed into
a path of shorter available phrases, with cost equal to the number of
pieces minus one.  It explores freedom in intermediate binary
associations beyond a fixed pool of single binary rules.  This is an
exploratory finite model; its endpoint orientation tests do not by
themselves certify an arbitrary completed binary realization.  At length
cap 16, it has 2,785 word variables, 473,057 flow edges, and 475,191
constraints.  HiGHS reached the 180-second time limit without an
incumbent, taking about 182.3 seconds in the measured call.  Its result is
\textbf{UNKNOWN}.  It provides no infeasibility evidence.

A separate dynamic-programming parser assigns low cost to phrases whose
dependencies are already paid for outside AX.  The unchanged non-AX
closure contains 51 helper rules, leaving 77 AX-exclusive rules in the
baseline.  The parser minimizes a weighted occurrence cost, then builds
an orientation-compatible binary association and counts its actual shared
closure.  Nine reuse weights were tried.  The control and 172 single
adjacent disjoint swaps gave 1,557 weighted parses; adding the 224
rank-generated words gave 3,573.  Both campaigns retained a best count of
128.  These are heuristic parses: minimizing a weighted occurrence cost
is not equivalent to minimizing a shared grammar globally.

Finally, a direct finite equality screen compared the 128 old grammar
words and 44 inherited atoms.  All 172 remain distinct under exact fine
words, disjoint-support traces, affine-extended traces, and a signature
that multiplies each maximal consecutive affine run while keeping
non-affine letters opaque.  This rules out an immediate merge by these
four equality tests on these existing objects.  Distinct signatures do
not prove distinct program semantics.  New factorizations or a stronger
semantic equality procedure remain possible.

\input{grammar/zero5_lock_obstruction.tex}

\subsection{Reproducing the recorded claims}

The finite replay uses only Python's standard library.  From the archive
root run
\begin{verbatim}
python3 grammar/reproduce.py verify
\end{verbatim}
It checks the local manifest, regenerates the word vocabularies and
search inputs from the pristine reference, rechecks all ten annealing
outputs and the MILP control, replays the four equality screens and the
five-edge obstruction, and recomputes the search summary from the saved
logs and outputs.  Its affine equality checks use the separate source
certificates described earlier.  It does not treat a solver's textual
status as a proof object.

Discovery runs are optional and distinct from this replay.  For example,
\begin{verbatim}
python3 grammar/reproduce.py search affine_trace
python3 grammar/reproduce.py search zero5_fully_unlocked
python3 grammar/reproduce.py milp fixed64
python3 grammar/reproduce.py parse ranked
\end{verbatim}
The C++ searches require a C++17 compiler; MILP discovery additionally
requires NumPy and SciPy.  Outputs go to an archive-relative
\file{grammar/reproduced/} directory by default, preserving the recorded
evidence.  The ten saved annealing budgets total about 42 minutes, before
compilation; a shorter optional run can use \texttt{--seconds}.  Full
parameter overrides and the inventory of code, source inputs, logs and
certificates are documented in \file{grammar/README.md}.
